% Sections 1--3 of the NAACL 2027 SRW submission.
% Fragment only: no preamble, no \begin{document}.
% Requires: acl-style class, booktabs, and the bibliography keys cited below.

\section{Introduction}
\label{sec:intro}

Chinese movie-review corpora are abundant in ratings and impoverished
in labels. Douban attaches a five-star rating to every short review,
but review-level semantic sentiment labels---positive, neutral, or
negative as expressed in the text itself---are essentially absent at
scale. The two signals are not interchangeable: a rating records an
overall evaluation, often recommendation intent, whereas the text may
praise one aspect and attack another. Reviews are also terse (median
24 characters), frequently ironic, and dense with film-specific
allusion; mixed evaluations make even ``neutral'' contestable
\citep{pang2005seeing,zhang2018dl}.

This paper studies what it takes to learn three-class sentiment
classification under this label scarcity, organized around three
research questions. \textbf{RQ1}: How reliably do user star ratings
represent review-level semantic sentiment, and how does this vary
across review populations? \textbf{RQ2}: How does LLM-assisted data
expansion affect weakly-supervised Chinese review sentiment
classifiers? \textbf{RQ3}: How do supervision sources and evaluation
protocols influence measured performance?

Our setting is a 910{,}998-review Douban corpus derived from
Moviedata-10M \citep{moviedata10m}, carrying star ratings but no
semantic labels, anchored by \emph{gold-600}: 600 human-verified
semantic labels (592 unique texts) scraped from a single film page.
On this base we build a weakly-supervised pipeline in the tradition
of self-training and programmatic weak supervision
\citep{lee2013pseudo,xie2020selftraining,ratner2017snorkel}: the gold
anchor trains a teacher ensemble that labels the full corpus,
confidence-filtered pseudo-labels train corpus-scale specialists, and
uncertainty-guided active learning adds $\sim$17K rows labelled by a
mixture of LLM teachers and human adjudication \citep{settles2009}.
Because every supervision source carries a different provenance and
failure mode, we evaluate with a source-aware protocol: per-row
provenance splits, a row-level train/test contamination audit, and a
decomposition of the leakage channels that inflate naive out-of-fold
estimates.

Our findings qualify each link of this chain. Star ratings agree with
human semantic labels on only 71.3\% of gold-600, and the gap is
population-dependent: on a 400-review identifier-disjoint blind test
agreement falls to 55.75\%---only 30.6\% of three-star reviews are
semantic-neutral on gold-600 (56.8\% on the blind sample)---and a
pure star-rating lookup already matches our deployed ensemble
(55.75\% vs.\ 55.8\%). LLM-assisted expansion does help: a 17K-label
mixture improves on the deployed ensemble by 9.0pp on the blind human
test, though LLM-consensus labels inflate measured accuracy by
roughly 26 points relative to human-adjudicated rows. Supervision-
source provenance moves conclusions in both directions: 68.7\% of our
development benchmark's rows were exposed to earlier corpus-derived
training, and transitive corpus exposure, second-order teacher
targets, and duplicate texts account for most of the naive-vs-corrected
OOF gap ($\sim$7pp three-class, $\sim$4pp binary; corrected figures
82.6--83.1\% / 91.5\%).

We make three contributions:
\begin{itemize}
  \item \textbf{An empirical analysis of the rating--sentiment gap}
        across three review populations (a curated single-film gold
        set, an identifier-disjoint blind human test, and the
        weakly-labelled corpus), locating where star supervision fails
        and how much a star-lookup baseline already explains.
  \item \textbf{A scalable weakly-supervised pipeline} combining
        star-derived weak labels, student pseudo-labels, and
        uncertainty-guided LLM annotation ($\sim$17K added labels),
        improving on the deployed ensemble by 9.0pp on
        identifier-disjoint human labels.
  \item \textbf{A source-aware evaluation protocol} for
        weakly-supervised pipelines---provenance-split reporting,
        row-level contamination audits, and leakage decomposition---
        that moves measured performance by several points.
\end{itemize}
Section~\ref{sec:related} reviews related work;
Section~\ref{sec:data} describes the datasets and annotation protocols;
Sections~\ref{sec:method}--\ref{sec:results} present the pipeline,
setup, and results; Section~\ref{sec:limitations} states limitations.

\section{Related Work}
\label{sec:related}

\subsection{Weakly Supervised Text Classification}

Learning from noisy or programmatic labels is mature: Snorkel
\citep{ratner2017snorkel} formalised weak supervision through
labelling functions; self-training \citep{lee2013pseudo} and its
noisy-student descendant \citep{xie2020selftraining}, extended by
consistency-training variants such as UDA \citep{xie2020uda}, show
that confident predictions can bootstrap larger training sets;
uncertainty-driven sampling underlies active learning
\citep{settles2009}; and loss correction \citep{arazo2021pseudo} and
confident learning \citep{northcutt2021} characterise which examples
are likely mislabeled. This literature treats label noise as a
property of the \emph{training} signal and reports on a fixed
evaluation split; we instead measure each supervision source's noise
structure directly (RQ1) and audit the evaluation itself (RQ3).

\subsection{LLM-Assisted Data Annotation}

LLM annotation is increasingly common: zero-shot LLMs can rival or
exceed crowd workers \citep{gilardi2023}, but two cautions apply.
First, LLM suggestions anchor human annotators toward model
predictions \citep{schroeder2025humanloop}; our adjudication
interface displayed the student model's prediction and confidence,
so our 479 human-adjudicated rows are explicitly exposed to this
effect (Section~\ref{sec:data}). Second, two passes of the
\emph{same} LLM share inductive bias rather than providing
independent annotation, and models trained on LLM-augmented data
inherit teacher conventions \citep{moller2024parrot}; on subjective
tasks annotator disagreement is informative rather than noise
\citep{davani2022}. We therefore treat LLM two-pass consensus as a
pseudo-label source, never ground truth \citep{zenimoto2024,neveditsin2025}.

\subsection{Review Sentiment Analysis}

Review sentiment analysis ranges from early work establishing that
star ratings and review polarity are related but distinct signals
\citep{pang2005seeing} through opinion-mining surveys
\citep{pang2008opinion} and deep models \citep{zhang2018dl}. Chinese
review sentiment relies on pre-trained encoders---MacBERT
\citep{cui2020macbert} and sentiment-enhanced SKEP
\citep{tian2020skep} are our backbone choices---and aspect-level
models for mixed evaluations \citep{wang2016attention}. Prior work
assumes text-level labels exist; we ask how far star ratings,
pseudo-labels, and LLM labels can substitute for them.

\subsection{Evaluation and Dataset Quality}

Data contamination in evaluation is increasingly documented for LLM
benchmarks, with taxonomies distinguishing incidental from targeted
exposure \citep{sainz2023contamination,balloccu2024}, and leakage is
a recognised reproducibility hazard \citep{kapoor2023}. For OOD
evaluation, classical detection baselines \citep{hendrycks2017} and
benchmarks such as \citet{baran2023ood} establish protocols we
adapt; reusing a development set for model selection introduces
measurable optimism \citep{cawley2010}; and
\citet{dror2018hitchhiker} provide the statistical-testing
discipline---McNemar tests, bootstrap intervals, Holm
correction---that we apply to every comparison. Reporting norms for
human evaluation and annotation \citep{kunilovskaya2026,alkhalifa2026}
and dataset documentation \citep{bender2018,gebru2021} motivate our
per-row provenance fields; unlike audits of released benchmarks
\citep{toossi2024}, we audit a pipeline we build ourselves.

\section{Dataset and Annotation}
\label{sec:data}

\subsection{Overview}

The task is three-class sentiment classification
$y \in \{\text{positive}, \text{neutral}, \text{negative}\}$ of Douban
movie short reviews. Corpus reviews are short (median 24 characters)
and frequently contain sarcasm, film-specific allusions, or mixed
polarity, so the neutral class is both frequent and ambiguous.
Table~\ref{tab:datasets} summarises every dataset, its annotation
source, and its role. Two properties are unusual and deliberate:
several ``test'' sets carry LLM-generated labels, and one of them
(ood1000) is a \emph{development/selection} set, not a held-out test.

\begin{table}[t]
\centering
\scriptsize
\setlength{\tabcolsep}{2pt}
\caption{Datasets and annotation sources. ``Consensus'' = agreement of
two blind LLM passes; ``adjudicated'' = human resolution of pass
disagreements plus spot-checked agreements. Roles reflect the audited,
corrected protocol. The $\sim$17K AL pools mix 279 human-adjudicated,
721 consensus, and $\sim$15.9K DeepSeek-V4.1-Flash teacher labels;
their text-level audit is in Appendix~\ref{app:human400}.}
\label{tab:datasets}
\begin{tabular}{@{}p{1.5cm}r p{2.9cm} p{1.6cm}@{}}
\toprule
Dataset & Rows & Annotation source & Role \\
\midrule
gold-600 & 600 & Human semantic labels; star-rule weak labels & Train anchor; OOF eval \\
pseudo pools & 910{,}998 & Star-rating rule; student pseudo-labels & Weak supervision \\
ood1000 & 1{,}000 & 800 LLM consensus + 200 human-adjudicated & Dev/selection set \\
al1000 & 1{,}000 & 721 LLM consensus + 279 human-adjudicated & AL training \\
AL expansions (al5000--al17k) & $\sim$17{,}000 & 279 human-adj.; 721 consensus; $\sim$15.9K DS-Flash teacher & AL training \\
fresh500 & 500 & DeepSeek-V4.1-Flash zero-shot labels & Dev replication check (partially exposed) \\
lab500 & 500 & Two LLM annotators; 331 consensus rows & Superseded pilot; partially exposed \\
human-400 & 400 & Blind single-annotator labels; length-stratified & Identifier-disjoint, single-annotator, length-stratified human test \\
\bottomrule
\end{tabular}
\end{table}

\subsection{Gold-600: human labels and the star-rating gap}

The gold set was scraped from the first 30 pages of a single Douban
movie page: 600 rows covering 592 unique reviews (316 positive / 92
neutral / 192 negative). Eight exact-duplicate text pairs were
\emph{not} removed before training or evaluation, so duplicates can
land in different cross-validation folds---a residual fold-isolation
violation analysed in Appendix~\ref{app:audit}. Each row carries two
labels: \texttt{label\_rule}, a star-derived three-class weak label,
and \texttt{label\_semantic}, a human-verified semantic label.
Fourteen rows lack a star rating and are excluded from rating-derived
supervision and the agreement analysis below; neutral-class decision
rules (Appendix~\ref{app:annotation}) are applied consistently across
gold, adjudication, and human-400 labelling. Reviews here are longer
than corpus reviews ($\sim$111 vs.\ 40 characters), and neutral rests
on only 92 examples. All 600 texts were sent to Alibaba PAI for a
baseline comparison; we disclose this egress in
Section~\ref{sec:limitations}.

The dual labels let us measure, rather than assume, star ratings as
sentiment supervision. On the 586 gold rows carrying both labels,
per-star agreement is 94.3\% (1$\star$), 85.9\% (2$\star$),
\textbf{30.6\%} (3$\star$), 78.7\% (4$\star$), and 96.9\% (5$\star$);
overall agreement is \textbf{71.3\%}. The failure is concentrated in
the middle band---51.3\% of three-star reviews are semantically
\emph{negative}, not neutral---so we use ratings only as weak
pretraining signal, never as gold labels.

\subsection{Weakly-labelled corpus}

The weak corpus derives from Moviedata-10M \citep{moviedata10m}: 4.43M
raw Douban comments cleaned and deduplicated to \textbf{910{,}998}
reviews spanning 20{,}966 films and 249{,}105 users, with star rating,
timestamp, and vote counts but no semantic labels. It is the source of
all pseudo-label pools: a student ensemble labels the full corpus, a
confidence-filtered subset (271{,}905 rows at $p \ge 0.7$) trains
corpus-scale specialists, and margin/entropy sampling selects the
active-learning pools below.

\subsection{LLM-assisted benchmarks: lab2000 and successors}

\paragraph{lab2000.} To measure---and then narrow---the
gold-to-corpus generalisation gap, we built \textbf{lab2000}: 2{,}000
corpus reviews in two disjoint subsets. \textbf{ood1000} samples
1{,}000 reviews uniformly at random (gold-deduplicated, cross-movie)
and served as the development benchmark for model selection.
\textbf{al1000} samples 1{,}000 reviews by margin/entropy sampling
from the deployed student's predictions---deliberately the hardest
rows---and is a training set, not a test set. Active-learning expansion
continued through al5000, al5000b, and disagreement-mined al3k rounds,
producing pools up to $\sim$17K rows of mixed provenance (279
human-adjudicated, 721 same-LLM-agent consensus, $\sim$15.9K
DeepSeek-V4.1-Flash teacher labels, plus a Qwen3-4B ablation); all AL
pools are verified text-disjoint from every evaluation set.

\paragraph{Annotation protocol.} Each lab2000 row was labelled twice by
two independent inference passes of the \emph{same} LLM annotator model
with different few-shot example sets, blind to star ratings and to each
other's output---single-model self-consistency, not two independent
annotators. Inter-pass agreement was 86.7\% ($\kappa=0.80$) on ood1000
and 79.8\% ($\kappa=0.67$) on al1000; the lower agreement on the
margin-selected set is itself evidence that AL mining selects
intrinsically ambiguous rows. All 335 inter-pass disagreements plus 144
spot-checked agreements were adjudicated by the author, yielding
\textbf{479 human-adjudicated rows} (200 in ood1000, 279 in al1000);
the remaining rows take two-pass consensus (Table~\ref{tab:datasets}).

\paragraph{fresh500 and lab500.} \textbf{fresh500}: 500 uniformly
sampled corpus reviews labelled zero-shot by DeepSeek-V4.1-Flash; our
audit finds zero text overlap with the gold/AL pools but 70.2\% of its
rows share a corpus identifier with earlier corpus-derived training,
so it is a replication check on a partially exposed distribution, not
a clean test. \textbf{lab500} (superseded pilot): 500 reviews labelled
by two LLM annotators with 331 consensus rows; a pilot CV script used
lab500 itself for training, so it is partially exposed.

\paragraph{human-400.} Our identifier-disjoint final test: 400 corpus
reviews sampled with length stratification (160/120/80/40 across
$\le$20, 21--50, 51--100, and $>$100 characters) after excluding
every previously used identifier pool, and blind-labelled by one
author under written guidelines (Appendix~\ref{app:annotation}).
Post-hoc audits verify zero comment-ID overlap with every training
pool of every evaluated system, and a text-level audit against
462{,}879 pool texts (Appendix~\ref{app:human400}) finds zero exact
or near-duplicate overlap plus only two ultra-short boilerplate
normalized matches in corpus pseudo-label pools. The
length-stratified sample is benchmark-specific, not a population
estimate; labels are positive 156 / neutral 179 / negative 65.
Star-implied labels agree with these human labels on only
\textbf{55.75\%} of rows (vs.\ 71.3\% on gold-600); among 5$\star$
reviews only 62.7\% are semantically positive. Single-annotator
labelling means no IAA is computable; we treat human-400 as our
primary OOD evidence while flagging the limitation in
Section~\ref{sec:limitations}.

\subsection{Annotation quality: what the agreement numbers do and do not show}

We report agreement following standard IAA practice
\citep{artstein2008} and the reporting norms of
\citeauthor{kunilovskaya2026} (\citeyear{kunilovskaya2026}) and
\citeauthor{alkhalifa2026} (\citeyear{alkhalifa2026}); dataset
documentation follows data-statement and datasheet conventions
\citep{bender2018,gebru2021} (Appendix~\ref{app:ethics}). Three
caveats qualify all agreement numbers below. First, a 200-row random
subset of ood1000 was blind-labelled by a second human annotator, but
because the released labels mix 47 human-adjudicated and 153
LLM-consensus rows, the measured 71.5\% raw / $\kappa=0.568$ (95\% CI
$[0.47,0.66]$) is a human-vs-\emph{reference} figure, not a pure
human--human IAA. Second, on the 47 adjudicated rows where both LLM
passes and the human label are available, the two passes agree on only
\textbf{29.8\%} ($\kappa=-0.18$) versus 86.7\% overall---consensus
labels are least reliable on exactly the hard rows that decide
rankings. Third, adjudication displayed the deployed student's
prediction and confidence, the anchoring setting of
\citeauthor{schroeder2025humanloop} (\citeyear{schroeder2025humanloop});
our 479 ``human'' labels may be biased toward stu3c-family outputs,
which matters because the same family produces the pseudo-labels being
evaluated; Sections \ref{sec:method}--\ref{sec:results} quantify each.
